\scanpage{94}
\mspage{93}
b) \deleted{\unsure{que}} $f$ est compatible avec opérations de $\mu$,
quand on fait opérer $\mu$ trivialement sur $C$ ---
\deleted{\unsure{c) $f$.}}

c) $\mu$ opère librement sur $\operatorname{Ob}C'$, et
$\operatorname{Ob}C'/\mu\longrightarrow\operatorname{Ob}C$
est \uline{bijectif} ($\Rightarrow f$ surjectif sur les objets)

d) \deleted{\unsure{Pourtant}} Condition à exprimer, sur la
structure des catégories fibres $C'_{\unsure{\xi}}$ ($\unsure{\xi}\in\operatorname{Ob}C$),
en tant que \inserted{$\pi$-}group. \deleted{stricts}\inserted{connexes}, avec gr.
d'opérateurs $\mu$.

Soit $\Phi$ \deleted{est} un $\pi$-groupoïde \inserted{connexe}, avec op.
de $\mu$ sur $C'$ (compatible avec $\pi$
structure), tel que $\operatorname{Ob}\Phi$ soit un
torseur sous $\mu$, soit $E$. Soient $x,y\in E$,
et considérons $\operatorname{Hom}_{\Phi}(x,y)$, c'est un
$\pi$-torseur, soit $H_{x,y}$, et on a, via l'op.
de $\mu$
\[
H_{x,y}\xrightarrow[\sim]{\alpha_*}H_{\alpha x,\alpha y}
\qquad\alpha\in\mu
\]
on trouve un système transitif d'iso. entre
les $H_{x,y}$ pour lesquels $y^{-1}x=u\in\mu$
est fixé, on peut les identifier entre
eux à un $H_u$. \deleted{\unsure{avec}} La composition
des homomorphismes s'exprime par des
isom canoniques
\[
H_u\mathbin{\wedge}H_v\xrightarrow{\sim}H_{uv}
\]
satisfaisant la condition d'associativité, \inserted{d'unité}, qu'il
